% Pokemon Red RL — EWRL 2026 paper draft
%
% Working title: "Symbols or Pixels? A Controlled Study of Observation
% Representations in Long-Horizon Reinforcement Learning"
%
% Build:    cd paper && make pdf
% Sections: paper/sections/0[1-6]_*.tex
% References: paper/references.bib
%
% NOTE: this uses a generic article class.  When the EWRL 2026 LaTeX
% template is published (typically 4-6 weeks before the deadline), swap
% \documentclass and any venue-specific commands.  Section content is
% template-agnostic.

\documentclass[11pt,a4paper]{article}

% ── Packages ──────────────────────────────────────────────────────────
\usepackage[utf8]{inputenc}
\usepackage[T1]{fontenc}
% microtype gives slightly tighter typesetting but requires font-expansion
% support that the TeX Live "basic" scheme doesn't always ship.  Use
% protrusion-only mode so it works on minimal TeX installs.
\usepackage[protrusion=true,expansion=false]{microtype}
\usepackage[margin=1in]{geometry}
\usepackage{amsmath,amssymb,amsthm}
\usepackage{graphicx}
\usepackage{booktabs}
\usepackage{xcolor}
\usepackage[hidelinks]{hyperref}
\usepackage{natbib}

% Lightweight cleveref-like macro: \Cref expands to "Section~\ref{...}"
% etc. Avoids the cleveref package which isn't in TeX Live basic.
\newcommand{\Cref}[1]{Section~\ref{#1}}
\newcommand{\cref}[1]{\Cref{#1}}

% ── Layout & style ────────────────────────────────────────────────────
\setlength{\parskip}{4pt}
\renewcommand{\baselinestretch}{1.05}

% Annotation macro for results-section placeholders.  Search for "TODO"
% to find every spot the pilot data needs to land.
\newcommand{\todopilot}[1]{\textcolor{red}{[TODO(pilots): #1]}}

% ── Title block ───────────────────────────────────────────────────────
\title{%
  Symbols or Pixels? \\
  \Large A Controlled Study of Observation Representations in Long-Horizon
  Reinforcement Learning
}
\author{%
  Alan Chester\thanks{CAM Labs LLC. Correspondence:
    \href{mailto:amcheste@gmail.com}{amcheste@gmail.com}}
}
\date{%
  EWRL 2026 submission \\[2pt]
  \small\today
}

\begin{document}

\maketitle

% ── Abstract ──────────────────────────────────────────────────────────
\begin{abstract}
We present a controlled empirical study of observation representations
for reinforcement-learning agents in a long-horizon, sparse-reward game
environment.  Using \emph{Pokémon Red} as a benchmark with the Boulder
Badge as the primary milestone, we compare three observation
treatments---\emph{pixel} (raw downsampled screen), \emph{symbolic}
(structured game state read from RAM), and \emph{hybrid} (both
streams)---under matched architecture, reward, and compute.
%
We train RecurrentPPO agents for 10M environment steps across three
seeds per treatment and report aggregate metrics with 95\% bootstrap
confidence intervals using \texttt{rliable}.
%
\todopilot{Insert one-sentence headline result: which treatment wins
sample efficiency, by how much, and whether the CI excludes the others.}
%
We release \textsc{PokeGym}, a Gymnasium-compatible benchmark, and the
full training and analysis pipeline under an open license to support
follow-up work.
\end{abstract}

% ── Body ──────────────────────────────────────────────────────────────
\section{Introduction}
\label{sec:intro}

Deep reinforcement-learning agents are typically benchmarked on
short-horizon tasks with rich, dense reward signals---most commonly the
Atari Learning Environment~\citep{bellemare2013arcade,mnih2015atari}
and procedurally-generated grid worlds~\citep{cobbe2020procgen,
matthews2024craftax,samvelyan2021minihack}.  These environments stress
exploration, sample efficiency, and generalisation, but few of them
require an agent to maintain coherent intent across the millions of
steps that real long-horizon games---role-playing games, strategy
games, or open-world games---demand.

\emph{Pok\'emon Red}, a 1996 Game Boy role-playing game, sits at an
unusually informative point in the difficulty landscape.  Its
state-action space is small enough to render with a Game Boy emulator
at thousands of frames per second on consumer hardware, but the game
contains long causal chains (deliver Oak's parcel before exiting Viridian;
beat Brock with a fire- or water-resistant team; navigate Mt. Moon
without losing the run) that punish agents which lack temporal
abstraction.  The 18 event flags between the player's bedroom and the
Boulder Badge encode the kind of structured progress that the dense
reward proxies of Atari obscure.

In this work we use Pok\'emon Red as a controlled test bed for one
specific question: \textbf{how does the choice of observation
representation affect what a long-horizon RL agent can learn?}  We
compare three treatments under matched architecture, reward, and
compute:
%
\begin{itemize}
\item \textbf{Pixel} --- the standard $80 \times 72$ grayscale screen
  observation, encoded with a Nature-DQN-style
  CNN~\citep{mnih2015atari}.
\item \textbf{Symbolic} --- a structured vector read directly from the
  emulator RAM, containing the local tile map, the 18 pre-registered
  event flags, party stats, and player coordinates, encoded with an
  MLP.
\item \textbf{Hybrid} --- both streams concatenated at the recurrent
  layer.
\end{itemize}
%
All three feed into a single LSTM and a PPO actor-critic
head~\citep{schulman2017ppo,hausknecht2015drqn}, with the same
event-flag-based reward function.  Differences in performance are
attributable to the observation representation alone.

\paragraph{Contributions.}
The contributions of this preliminary study, intended to seed a longer
journal-length follow-up, are:
\begin{enumerate}
\item \textbf{An open benchmark.}  We release \textsc{PokeGym}, a
  Gymnasium-compatible~\citep{gymnasium} environment for Pok\'emon Red
  with three first-class observation treatments, an event-flag reward
  calculator, and a curriculum-aware save-state system.  The
  environment runs headlessly at $\sim$\todopilot{N}~steps/s/CPU on a
  consumer laptop.
\item \textbf{A pre-registered analysis plan.}  Hypotheses, primary
  metric (Boulder Badge win-rate), seed counts, statistical tests
  (\texttt{rliable} stratified bootstrap with 95\% CIs), and stopping
  rules were committed to version control before the main experiments
  ran.\footnote{See \texttt{paper/analysis\_plan.md} in the open
    repository.}
\item \textbf{Empirical comparison of three observation treatments at
  10M~environment steps each, three seeds per treatment.}
  \todopilot{Replace this bullet once results are in: 1-2 sentences on
    the dominant treatment and the magnitude of the gap, with the
    bootstrap CI.}
\end{enumerate}

\paragraph{Scope.}
The pilot results reported here use 10M environment steps per seed and
three seeds per treatment.  We treat them as preliminary evidence
rather than a final ranking---both the seed count and the training
budget will be substantially expanded in subsequent journal-length
work, as committed in the pre-registered plan.

The remainder of the paper is organised as follows.
\Cref{sec:related-work} situates the work in the long-horizon-RL,
representation-learning, and Pok\'emon-RL literatures.
\Cref{sec:environment} describes \textsc{PokeGym} and the three
observation treatments in detail.  \Cref{sec:methods} fixes the
training pipeline, statistical methodology, and pre-registered
deviations.  \Cref{sec:results} presents pilot learning curves,
aggregate metrics, and milestone progression.  \Cref{sec:discussion}
discusses what the pilot can and cannot conclude, and outlines the
follow-up experiments planned for the journal version.

\section{Related Work}
\label{sec:related-work}

\paragraph{Long-horizon RL benchmarks.}
The Arcade Learning Environment~\citep{bellemare2013arcade} provided
the dominant benchmark for value-based and policy-gradient deep RL
methods through the late 2010s, but its tasks rarely exceed a few
thousand frames per episode and have dense, well-shaped reward
signals.  Procgen~\citep{cobbe2020procgen} added procedural generation
to stress generalisation but kept per-episode horizons short.  More
recent open-ended benchmarks---Craftax~\citep{matthews2024craftax} and
MiniHack~\citep{samvelyan2021minihack}---move toward the
sparse-reward, long-horizon regime, but do so in synthetic worlds whose
dynamics are simple enough to render in compute-light JAX.  The
present work targets the same regime via a real, hand-designed game
environment, where the causal chains and reward sparsity are products
of human design rather than procedural generation.

\paragraph{Pok\'emon Red as an RL test bed.}
The Pok\'emon-RL community has produced influential code releases and
blog-post-quality empirical results on Pok\'emon
Red~\citep{whitney2023pokemonrl,pokemonred_experiments}, demonstrating
that PPO with a pixel-based observation can learn meaningful behaviour
within the first few towns of the game.  These projects have been
crucial for showing the environment is tractable, but they do not
present controlled comparisons across observation treatments, do not
report bootstrap-based aggregate statistics, and have not been
released as installable Gymnasium environments.  We build directly on
their foundation while adding the methodological scaffolding required
for publication-ready empirical claims.

\paragraph{Observation representations in RL.}
The pixel-vs-state question has been studied in narrower contexts:
\citet{anand2019unsupervised} showed that learned state representations
on Atari preserve much of the underlying game state, suggesting that
pixel agents \emph{re-derive} information that a symbolic agent could
read directly.  \citet{kornblith2019cka} introduced CKA as a tool for
comparing internal representations across networks, which we plan to
apply in follow-up work to ask whether pixel and symbolic agents
converge to similar internal features.  To our knowledge, no prior
work performs a matched-architecture, matched-reward, matched-compute
comparison of pixel, symbolic, and hybrid observation streams in a
single long-horizon RL task.

\paragraph{Recurrent policies under partial observability.}
Pok\'emon Red is partially observable from any single frame: the
overworld map, the player's party, and the event-flag state of the
world are not visible simultaneously.  Recurrent policies are the
canonical response to partial
observability~\citep{hausknecht2015drqn}.  We use the
RecurrentPPO~\citep{schulman2017ppo} implementation from
\texttt{sb3-contrib}~\citep{sb3} for all three treatments, holding the
recurrent-layer architecture fixed and varying only the upstream
encoder.

\paragraph{Statistical reporting in deep RL.}
\citet{agarwal2021rliable} catalogued the field's practice of reporting
single-seed learning curves and showed how strongly that practice
overstates effect sizes.  Their \texttt{rliable} library has become
standard for reporting interquartile means with stratified-bootstrap
confidence intervals, and we adopt it throughout.  Our pre-registered
plan commits to IQM as the primary point estimate and 95\% bootstrap
CIs with 2,000 resamples as the primary uncertainty estimate.

\section{Environment}
\label{sec:environment}

We release \textsc{PokeGym}, a Gymnasium-compatible
environment~\citep{gymnasium} that wraps a headless
PyBoy~\citep{pyboy} emulator running an unmodified Pok\'emon Red ROM.
The environment exposes a discrete action space, three first-class
observation representations, an event-flag-based reward calculator, and
a curriculum-aware save-state system.

\subsection{Action Space}
The base environment supports the standard Game Boy buttons: D-Pad
(\textsc{Up}, \textsc{Down}, \textsc{Left}, \textsc{Right}),
\textsc{A}, \textsc{B}, \textsc{Start}.  We exclude \textsc{Select}
because it has no in-game effect for the section of the game we
study.  This yields $|\mathcal{A}| = 7$.  Each environment step
advances the emulator by 24 frames (one in-game tile of overworld
movement) before returning control to the agent.

\subsection{Observation Treatments}
\label{sec:obs-treatments}

The same emulator state can be exposed as three different observation
spaces, selected at environment construction time via
\texttt{observation\_type}.

\paragraph{Pixel.}
A single $80 \times 72$ grayscale frame ($\mathcal{B}^{80 \times 72}$,
$\mathcal{B} = \{0, \ldots, 255\}$).  We use grayscale rather than RGB
because the Game Boy's 4-shade palette already loses no information
relative to RGB, and this halves the input dimensionality.  The pixel
agent encodes this with a Nature-DQN-style
CNN~\citep{mnih2015atari}---three convolutional layers followed by a
fully-connected projection to a 256-dim feature vector.

\paragraph{Symbolic.}
A flat 29-dimensional float vector containing:
%
\begin{itemize}
\item \textbf{Player position} (3 dims): map ID, $x$ tile coordinate,
  $y$ tile coordinate.
\item \textbf{Health and party} (4 dims): current HP fraction, party
  size, max party level, total levels summed.
\item \textbf{Currency and inventory} (3 dims): money (log-scaled),
  badge count, item-bag fill ratio.
\item \textbf{Exploration state} (2 dims): unique-maps-visited count,
  unique-tiles-visited count.
\item \textbf{Event flags} (17 dims): a binary vector for the 17
  pre-registered Boulder Path event flags (\Cref{sec:events}); the
  18th flag (\texttt{EVENT\_GOT\_POKEMON\_FROM\_FAN\_CLUB\_CHAIRMAN})
  is held out as a sanity-check canary.
\end{itemize}
%
The symbolic agent encodes this with a 2-layer MLP ($29 \to 128 \to
256$) projecting to the same 256-dim feature space as the pixel CNN,
keeping the downstream LSTM input dimension matched.

\paragraph{Hybrid.}
A \texttt{Dict} space combining both branches.  The CNN and MLP run in
parallel; their 256-dim outputs are concatenated to a 512-dim vector
before entering the LSTM.  The hybrid agent has roughly 2$\times$ the
encoder parameters of pixel or symbolic alone, but identical
recurrent and policy heads.  We do not match parameter counts across
treatments---matching the encoder \emph{capacity} would require
shrinking the pixel CNN below the published Nature-DQN architecture
and would conflate two design decisions; we match the architecture
above the encoder instead.  We discuss this trade-off in
\Cref{sec:discussion}.

\subsection{Reward}
\label{sec:reward}

We use an event-flag-based reward calculator that issues a fixed
positive reward the first time each pre-registered flag is triggered
in an episode, plus a small per-step time penalty.  The 18-flag set is
committed in the pre-registered plan and corresponds to the canonical
critical path between the player's bedroom in Pallet Town and Brock's
Boulder Badge in Pewter City Gym.  In addition to event flags, we
issue:
%
\begin{itemize}
\item A small per-step time penalty ($-0.001$) to encourage timely
  progress.
\item A discovery reward for the first visit to each new map ID.
\item Negative reward for fainting (HP $\to$ 0 across the entire
  party).
\end{itemize}
%
All other shaping signals (e.g., level-based or money-based shaping)
are zero by default.  This keeps the reward signal close to the
event-flag milestones, which are observable, interpretable, and
directly map to the game's progression.

\subsection{Pre-Registered Event Flags}
\label{sec:events}

The 18 event flags used by the reward calculator are listed in the
appendix and committed in the pre-registered plan.  They cover:
deliver Oak's parcel (8 flags), traverse Viridian Forest (3 flags),
defeat Brock's gym trainers and Brock himself (5 flags), plus one
optional stretch milestone (Rival fight on Route 22) and one canary
flag that should not be reachable before Brock and is monitored as a
sanity check.

Flag bits and addresses are extracted from the
\texttt{pret/pokered}\footnote{\url{https://github.com/pret/pokered}}
disassembly and verified by emulator-side replay.

\subsection{Curriculum and Episode Resets}
\label{sec:curriculum}

We use a single fixed save state, \texttt{post\_intro.state}, captured
in Pallet Town immediately after Professor Oak's introduction
sequence.  All episodes reset to this state.  Pilot results in this
paper therefore reflect the agent's ability to traverse Route 1 →
Viridian → Viridian Forest → Pewter from a fixed starting point.
Future work will sample from a curriculum of save states (Pallet,
Viridian Pok\'emon Center, Mt. Moon entrance) using a Go-Explore-style
weighted distribution.

\subsection{Episode Truncation}
Episodes truncate at 15{,}000 environment steps regardless of in-game
state.  This is approximately 30 minutes of in-game time at the
emulator's 24-frame action repeat and is empirically sufficient for
even a randomly-acting policy to leave Pallet Town occasionally,
preventing the agent from being stuck in a state with zero reward
signal.

\subsection{Reproducibility}
\textsc{PokeGym} is open-source and pip-installable.  The training
script, evaluation harness, statistical analysis pipeline, and live
training dashboards are released under an OSI-approved license at the
project repository.  We seed every RNG that affects training (Python,
NumPy, PyTorch, Gymnasium) deterministically per run and document the
exact seeds in the supplementary material.

\section{Methods}
\label{sec:methods}

This section fixes the training pipeline and the statistical
methodology.  The values reported here were committed to the
pre-registered plan in \texttt{paper/analysis\_plan.md} before any
main-result experiments were run; we note any deviations explicitly in
\Cref{sec:deviations}.

\subsection{Algorithm and Architecture}
\label{sec:algorithm}

We use RecurrentPPO from
\texttt{sb3-contrib}~\citep{sb3,schulman2017ppo,hausknecht2015drqn} for
all three treatments.  The architecture upstream of the LSTM differs
across treatments (\Cref{sec:obs-treatments}) but everything from the
LSTM onward is identical:
%
\begin{itemize}
\item \textbf{LSTM:} 1 layer, 256 hidden units (parameter-matched
  across all three treatments).
\item \textbf{Policy and value heads:} 2-layer MLPs of width 64 each,
  with \texttt{tanh} activations.
\item \textbf{Action distribution:} categorical over the 7 buttons.
\end{itemize}
%
The decision to vary only the encoder architecture is deliberate: any
performance difference between treatments is attributable to what the
observation stream makes available to the LSTM, not to differences in
recurrence or policy expressivity.

\subsection{Hyperparameters}
\label{sec:hyperparameters}

The RecurrentPPO hyperparameters used in the pilot are:
\todopilot{Confirm against the actual values in
\texttt{pokemon\_red\_ai/training/models.py} after the first run; the
values below are the pre-registered defaults.}
%
\begin{table}[h]
\centering
\small
\begin{tabular}{lr}
\toprule
Hyperparameter & Value \\
\midrule
Learning rate & $3 \times 10^{-4}$ \\
Batch size    & $256$ \\
Rollout length ($n\_\text{steps}$) & $2{,}048$ \\
Epochs per update & $10$ \\
Entropy coefficient & $0.01$ \\
Discount $\gamma$ & $0.99$ \\
GAE $\lambda$ & $0.95$ \\
Clip range & $0.2$ \\
Parallel envs & $1$ \\
\bottomrule
\end{tabular}
\end{table}
%
We deliberately use a single environment per training run rather than
the 8-env parallelism committed in the long-form plan, because the
PyBoy emulator is single-threaded per instance and parallelism on a
single machine yields diminishing returns past 4 envs in our
preliminary timing.  For the journal version we will revisit
parallelism once compute scales beyond a single workstation.

\subsection{Training Protocol}
\label{sec:protocol}

Each treatment is trained for $10^7$ environment steps across three
seeds (42, 123, 456).  No hyperparameter tuning is performed within
the pilot; all treatments use identical hyperparameters.  Training is
checkpoint-saved every 250{,}000 steps and the best checkpoint per run
(by training-time best-episode reward) is retained for evaluation.

\subsection{Evaluation Metric and Statistical Analysis}
\label{sec:stats}

The pre-registered \emph{primary} metric is Boulder Badge win-rate at
the end of training, evaluated as the fraction of 20 deterministic
rollouts that earn the badge before episode truncation.  At
10M~environment steps---one-tenth of the journal-version
budget---reaching Brock is a stretch goal rather than a baseline
expectation; in this paper we therefore report the primary metric
together with two informative secondary metrics:
%
\begin{itemize}
\item \textbf{Best episode reward} during training (accumulated event
  flags + discovery + time penalty).  This is a continuous proxy for
  progress and is what we use as the single number per seed when
  computing IQM.
\item \textbf{First episode at which each event flag is triggered.}
  Reported per-treatment as a milestone progression chart.
\end{itemize}

\paragraph{Aggregate metrics with confidence intervals.}
We use \texttt{rliable}~\citep{agarwal2021rliable} for all aggregate
statistics.  The primary point estimate per treatment is the
\emph{interquartile mean} (IQM) over the seed-level best-episode
rewards; we drop the bottom and top quartiles before averaging to
reduce sensitivity to outlier seeds.  Uncertainty is reported as a
95\% percentile bootstrap CI with $2{,}000$ resamples.

\paragraph{Probability of improvement.}
For pairwise comparisons we report the probability of improvement
(PoI), $\Pr[\text{score}_A > \text{score}_B]$, estimated via stratified
bootstrap.  We treat a treatment as superior to another only if the
95\% CI of PoI excludes 0.5.

\subsection{Deviations from the Pre-Registered Plan}
\label{sec:deviations}

The following deviations from the pre-registered plan are made
explicitly for the EWRL pilot and will be returned to nominal for the
journal version:
%
\begin{itemize}
\item \textbf{Seed count.}  Three seeds rather than the
  five-minimum committed for the main result.  We adopt three for the
  pilot to fit within the EWRL submission window; the journal version
  will use five (minimum) or eight (preferred) seeds per treatment.
\item \textbf{Training budget.}  $10^7$ steps rather than $10^8$.
  Reaching Brock is a stretch goal at this budget.  Effects detectable
  at $10^7$ are likely conservative estimates of effects at the
  full-budget version.
\item \textbf{Single environment per run.}  We run one parallel
  environment rather than the eight committed in the plan.  The
  journal version will revisit this once compute scales beyond a
  single workstation.
\end{itemize}
%
We do not deviate from the pre-registered hypotheses, primary or
secondary metrics, statistical tests, event-flag set, or evaluation
protocol.

\subsection{Compute and Reproducibility}
\label{sec:compute}

\todopilot{Fill from \texttt{paper/compute\_ledger.md} after pilots
finish: per-treatment wall-clock time, GPU utilisation peak, peak
memory, total GPU-hours.  Format as a small table.}

The training script, the evaluation harness, the statistical
pipeline, and the live training dashboards are open-sourced in the
project repository.  Training-run hyperparameter dumps and W\&B
artifact links for the runs reported here are provided in the
supplementary material.

\section{Results}
\label{sec:results}

% NOTE: This section is a structural skeleton.  Every \todopilot{...}
% marker is a place where pilot data needs to be inserted once the
% nine training runs finish.  See docs/research_playbook.md for the
% commands that produce these data; figures live in paper/figures/.

\subsection{Setup recap}

Three observation treatments---\emph{pixel}, \emph{symbolic},
\emph{hybrid}---trained for $10^7$ environment steps each, three seeds
per treatment, identical hyperparameters
(\Cref{sec:hyperparameters}).  All other components of the training
pipeline are matched (\Cref{sec:methods}).

\subsection{Learning Curves}
\label{sec:results-curves}

\begin{figure}[t]
\centering
% Produced by scripts/analyze.py and scripts/compare.py
% \includegraphics[width=0.85\linewidth]{figures/learning_curves.pdf}
\caption{%
  \todopilot{Learning curves: per-episode total reward smoothed with a
    20-episode rolling mean, with mean ($\pm$ std) bands across the
    three seeds for each treatment.  Replace this caption with one or
    two sentences describing the qualitative shape: which treatment
    rises fastest, which plateaus, whether bands overlap.}
}
\label{fig:learning-curves}
\end{figure}

\todopilot{Insert one paragraph describing \Cref{fig:learning-curves}
qualitatively.  Mention: the treatment with the steepest rise, the
treatment with the highest asymptote within the 10M-step budget,
whether the curves cross, and whether the across-seed std bands
overlap meaningfully.}

\subsection{Aggregate Metrics}
\label{sec:results-aggregate}

\begin{table}[t]
\centering
\small
\caption{%
  Aggregate metrics per treatment.  IQM is the interquartile mean of
  per-seed best-episode reward; CIs are 95\% percentile bootstrap
  intervals with 2{,}000 resamples.  $n$ = 3 seeds per treatment.
}
\label{tab:aggregate}
\begin{tabular}{lcccc}
\toprule
Treatment & Mean & Median & IQM & 95\% CI \\
\midrule
Pixel    & \todopilot{x.x} & \todopilot{x.x} & \todopilot{x.x} & \todopilot{[lo, hi]} \\
Symbolic & \todopilot{x.x} & \todopilot{x.x} & \todopilot{x.x} & \todopilot{[lo, hi]} \\
Hybrid   & \todopilot{x.x} & \todopilot{x.x} & \todopilot{x.x} & \todopilot{[lo, hi]} \\
\bottomrule
\end{tabular}
\end{table}

\todopilot{Insert one paragraph interpreting \Cref{tab:aggregate}.
Highlight: the highest-IQM treatment, whether its CI overlaps the
others, and the relative spread of the bands across seeds (e.g., "the
hybrid treatment had a tighter CI, suggesting more consistent
training across seeds").}

\paragraph{Probability of improvement.}
\todopilot{Insert pairwise PoI numbers from
\texttt{scripts/analyze.py}'s probability-of-improvement matrix.  For
each ordered pair (A, B), report $\Pr[\text{score}_A >
\text{score}_B]$ with its 95\% CI.  Mention which pairs cross or fail
to cross 0.5.}

\subsection{Milestone Progression}
\label{sec:results-milestones}

\begin{figure}[t]
\centering
% \includegraphics[width=0.95\linewidth]{figures/milestone_race.pdf}
\caption{%
  \todopilot{First-episode race chart for the 18 pre-registered event
    flags.  For each treatment, the median first-triggered episode
    across seeds is plotted; flags never triggered by any seed are
    shown as a hatched bar at the right edge.  Lower is faster.}
}
\label{fig:milestones}
\end{figure}

\todopilot{Insert one paragraph reporting which milestones each
treatment reached.  Pay particular attention to: (i) the
\texttt{EVENT\_GOT\_STARTER} flag, which any minimally-functional
agent should reach quickly; (ii) the \texttt{EVENT\_BEAT\_BROCK} flag,
which is the primary metric and is unlikely to be reached at 10M
steps; (iii) the canary flag
\texttt{EVENT\_GOT\_POKEMON\_FROM\_FAN\_CLUB\_CHAIRMAN}, which should
\emph{not} be reached---confirming that the agent is not exploiting an
unintended path.}

\subsection{Final-Window Performance}
\label{sec:results-final}

\begin{figure}[t]
\centering
% \includegraphics[width=0.6\linewidth]{figures/final_performance.pdf}
\caption{%
  Final-window performance: mean reward over the last 50 episodes per
  seed, aggregated as mean $\pm$ std across the three seeds per
  treatment.
}
\label{fig:final-perf}
\end{figure}

\todopilot{One short paragraph: agreement or disagreement between the
final-window bar chart and the IQM table.  Disagreement would be
informative---e.g., a treatment with high IQM but low final-window
reward suggests instability late in training.}

\subsection{Compute}
\label{sec:results-compute}

\todopilot{Add a small table with wall-clock per run (mean across
seeds), GPU model, peak GPU utilisation, and total GPU-hours for the
nine pilot runs.  Pull from \texttt{paper/compute\_ledger.md}.}

\section{Discussion}
\label{sec:discussion}

\subsection{What the Pilot Can and Cannot Conclude}

\todopilot{Open with one paragraph summarising the headline claim:
e.g., ``At 10M environment steps and three seeds per treatment, the
\emph{X} observation treatment achieved an IQM of \emph{Y}, with a
95\% bootstrap CI \emph{Z}; under the pre-registered probability-of-
improvement threshold, this difference is/is not statistically
significant.''  Be precise about what is and is not concluded.}

The 10M-step budget and three-seed pilot is deliberately preliminary.
It is sufficient to detect large effect sizes---a treatment that
clearly cannot leave Pallet Town versus one that consistently reaches
Viridian Forest---but it cannot reliably distinguish small-to-medium
effects.  In particular, the bootstrap CIs for the three treatments
\todopilot{do/do not} overlap, and we therefore \todopilot{can/cannot}
make a strong claim about the ranking of treatments at the
journal-version compute budget.

\subsection{Hypothesis 1 Revisited}
\label{sec:h1}

The pre-registered first hypothesis (H1) was that \emph{at matched
compute and matched architecture, symbolic observations yield
significantly better sample efficiency than pixel observations}.

\todopilot{Address H1 directly.  Did the data support, refute, or
fail-to-distinguish the hypothesis?  Mention the specific comparison
(IQM-symbolic vs IQM-pixel + their CIs) and the corresponding
probability-of-improvement value.  Be honest if the pilot is
inconclusive.}

\subsection{Why Pixel Might Underperform}

A pixel agent must implicitly re-derive all the structured state that
the symbolic agent reads directly from
RAM~\citep{anand2019unsupervised}: the player's coordinates, the
event-flag bit-vector, the contents of its party.  At a fixed compute
budget that re-derivation cost is paid in policy-network capacity and
in sample efficiency.  In a long-horizon, sparse-reward environment
where reward is gated on event flags, the symbolic agent has direct
access to the very signal that determines reward.  We expect this gap
to be especially pronounced early in training, when the pixel agent's
internal representations are not yet useful.

\subsection{Why Hybrid Might Not Strictly Dominate}

A naïve expectation is that the hybrid treatment, having access to
\emph{both} streams, should always win.  In practice, with a fixed
training budget, the hybrid agent must learn to weight two encoders
and may suffer from gradient interference if the two streams provide
redundant or conflicting signals.  Whether hybrid actually wins, ties,
or loses against the better single-stream treatment is an empirical
question; \todopilot{report the observed ranking}.

\subsection{Threats to Validity}

\paragraph{Encoder capacity is not parameter-matched.}
The hybrid encoder has roughly twice the parameter count of pixel or
symbolic alone.  This is a known confound; we discuss the trade-off in
\Cref{sec:obs-treatments}.  A capacity-matched ablation, in which the
single-stream encoders are widened to match the hybrid encoder's
total parameter count, is part of the journal-version plan.

\paragraph{Single fixed save state.}
All training in this pilot resets to \texttt{post\_intro.state}
(Pallet Town).  Effects observed here may not generalise to a
curriculum that samples from multiple save states, e.g., one rooted
inside a more partially-observable region of the game such as Mt.
Moon.  We will address this in the journal version with a
Go-Explore-style weighted save-state curriculum.

\paragraph{Three seeds.}
Three seeds is enough to detect large effects but provides limited
power for medium effects, particularly at the IQM tails.  All
quantitative claims are reported with bootstrap CIs and a
probability-of-improvement threshold; we do not over-interpret point
estimates.

\paragraph{No hyperparameter tuning per treatment.}
We use the same RecurrentPPO hyperparameters across all three
treatments, which is fair but pessimistic for a treatment whose
optimum lies far from those defaults.  The journal version will
include per-treatment hyperparameter tuning via Optuna on a held-out
save state.

\subsection{Future Work}
\label{sec:future-work}

The journal-length follow-up will extend this pilot in three
directions.  First, the seed count will increase to five (minimum)
and the training budget to $10^8$ steps per seed, restoring the
pre-registered design.  Second, four pre-committed
ablations---removing the LSTM, removing RND-style intrinsic
exploration, ablating frame-stacking, and ablating the save-state
curriculum---will isolate the contributions of recurrence, novelty
bonuses, temporal observation, and curriculum to the per-treatment
gap.  Third, a representational similarity analysis using
CKA~\citep{kornblith2019cka} will ask whether the three treatments
converge to similar internal features once trained, or whether the
performance gap reflects a deeper structural difference in what each
observation stream lets the agent represent.

\subsection{Open Release}

The full training, evaluation, and analysis pipeline is open-sourced.
The environment is pip-installable as \textsc{PokeGym}.  The W\&B
project is public and the run-level logs and model checkpoints are
linked from the supplementary material.  We hope the benchmark and
the analysis scaffolding lower the barrier for further work in
long-horizon RL on real, hand-designed game environments.


% ── References ────────────────────────────────────────────────────────
\bibliographystyle{plainnat}
\bibliography{references}

\end{document}
